\section{Quasi-regular sequences}
\label{section-quasi-regular}

\noindent
We introduce the notion of quasi-regular sequence which is slightly weaker
than that of a regular sequence and easier to use.
Let $R$ be a ring and let $f_1, \ldots, f_c \in R$.
Set $J = (f_1, \ldots, f_c)$. Let $M$ be an $R$-module.
Then there is a canonical map
\begin{equation}
\label{equation-quasi-regular}
M/JM \otimes_{R/J} R/J[X_1, \ldots, X_c]
\longrightarrow
\bigoplus\nolimits_{n \geq 0} J^nM/J^{n + 1}M
\end{equation}
of graded $R/J[X_1, \ldots, X_c]$-modules defined by the rule
$$
\overline{m} \otimes X_1^{e_1} \ldots X_c^{e_c} \longmapsto
f_1^{e_1} \ldots f_c^{e_c} m \bmod J^{e_1 + \ldots + e_c + 1}M.
$$
Note that (\ref{equation-quasi-regular}) is always surjective.

\begin{definition}
\label{definition-quasi-regular-sequence}
Let $R$ be a ring.
Let $M$ be an $R$-module.
A sequence of elements $f_1, \ldots, f_c$ of $R$ is called
{\it $M$-quasi-regular} if (\ref{equation-quasi-regular})
is an isomorphism. If $M = R$, we call $f_1, \ldots, f_c$ simply a
{\it quasi-regular sequence}.
\end{definition}

\noindent
So if $f_1, \ldots, f_c$ is a quasi-regular sequence, then
$$
R/J[X_1, \ldots, X_c] = \bigoplus\nolimits_{n \geq 0} J^n/J^{n + 1}
$$
where $J = (f_1, \ldots, f_c)$. It is clear that being a quasi-regular
sequence is independent of the order of $f_1, \ldots, f_c$.

\begin{lemma}
\label{lemma-regular-quasi-regular}
Let $R$ be a ring.
\begin{enumerate}
\item A regular sequence $f_1, \ldots, f_c$ of $R$ is a quasi-regular
sequence.
\item Suppose that $M$ is an $R$-module and that $f_1, \ldots, f_c$
is an $M$-regular sequence. Then $f_1, \ldots, f_c$ is an
$M$-quasi-regular sequence.
\end{enumerate}
\end{lemma}

\begin{proof}
Set $J = (f_1, \ldots, f_c)$.
We prove the first assertion by induction on $c$.
We have to show that given any relation
$\sum_{|I| = n} a_I f^I \in J^{n + 1}$ with $a_I \in R$ we
actually have $a_I \in J$ for all multi-indices $I$. Since
any element of $J^{n + 1}$ is of the form $\sum_{|I| = n} b_I f^I$
with $b_I \in J$ we may assume, after replacing $a_I$ by $a_I - b_I$,
the relation reads $\sum_{|I| = n} a_I f^I = 0$. We can rewrite
this as
$$
\sum\nolimits_{e = 0}^n
\left(
\sum\nolimits_{|I'| = n - e}
a_{I', e} f^{I'}
\right)
f_c^e
=
0
$$
Here and below the ``primed'' multi-indices $I'$ are required to be of the form
$I' = (i_1, \ldots, i_{c - 1}, 0)$. We will show by
induction on $l \in \{0, \ldots, n\}$
that if we have a relation
$$
\sum\nolimits_{e = 0}^l
\left(
\sum\nolimits_{|I'| = n - e}
a_{I', e} f^{I'}
\right)
f_c^e
=
0
$$
then $a_{I', e} \in J$ for all $I', e$.
Namely, set $J' = (f_1, \ldots, f_{c-1})$.
Observe that $\sum\nolimits_{|I'| = n - l} a_{I', l} f^{I'}$
is mapped into $(J')^{n - l + 1}$ by $f_c^{l}$.
By induction hypothesis (for the induction on $c$)
we see that $f_c^l a_{I', l} \in J'$.
Because $f_c$ is not a zerodivisor on $R/J'$ (as $f_1, \ldots, f_c$
is a regular sequence) we conclude that $a_{I', l} \in J'$.
This allows us to rewrite the term
$(\sum\nolimits_{|I'| = n - l} a_{I', l} f^{I'})f_c^l$
in the form $(\sum\nolimits_{|I'| = n - l + 1} f_c b_{I', l - 1}
f^{I'})f_c^{l-1}$. This gives a new relation of the form
$$
\left(\sum\nolimits_{|I'| = n - l + 1}
(a_{I', l-1} + f_c b_{I', l - 1}) f^{I'}\right)f_c^{l-1}
+
\sum\nolimits_{e = 0}^{l - 2}
\left(
\sum\nolimits_{|I'| = n - e}
a_{I', e} f^{I'}
\right)
f_c^e
=
0
$$
Now by the induction hypothesis (on $l$ this time) we see that
all $a_{I', l-1} + f_c b_{I', l - 1} \in J$ and
all $a_{I', e} \in J$ for $e \leq l - 2$. This, combined with
$a_{I', l} \in J' \subset J$ seen above, finishes the proof of the
induction step.

\medskip\noindent
The second assertion means that given any formal expression
$F = \sum_{|I| = n} m_I X^I$, $m_I \in M$ with $\sum m_I f^I
\in J^{n + 1}M$, then all the coefficients $m_I$ are in $J$.
This is proved in exactly the same way as we prove the corresponding
result for the first assertion above.
\end{proof}

\begin{lemma}
\label{lemma-flat-base-change-quasi-regular}
Let $R \to R'$ be a flat ring map. Let $M$ be an $R$-module.
Suppose that $f_1, \ldots, f_r \in R$ form an $M$-quasi-regular sequence.
Then the images of $f_1, \ldots, f_r$ in
$R'$ form a $M \otimes_R R'$-quasi-regular sequence.
\end{lemma}

\begin{proof}
Set $J = (f_1, \ldots, f_r)$, $J' = JR'$ and $M' = M \otimes_R R'$.
We have to show the canonical map
$\mu : R'/J'[X_1, \ldots X_r] \otimes_{R'/J'} M'/J'M' \to
\bigoplus (J')^nM'/(J')^{n + 1}M'$ is an isomorphism.
Because $R \to R'$ is flat the sequences
$0 \to J^nM \to M$ and
$0 \to J^{n + 1}M \to J^nM \to J^nM/J^{n + 1}M \to 0$
remain exact on tensoring with $R'$. This first implies that
$J^nM \otimes_R R' = (J')^nM'$ and then that
$(J')^nM'/(J')^{n + 1}M' = J^nM/J^{n + 1}M \otimes_R R'$.
Thus $\mu$ is the tensor product of (\ref{equation-quasi-regular}),
which is an isomorphism by assumption,
with $\text{id}_{R'}$ and we conclude.
\end{proof}

\begin{lemma}
\label{lemma-quasi-regular-sequence-in-neighbourhood}
Let $R$ be a Noetherian ring. Let $M$ be a finite $R$-module.
Let $\mathfrak p$ be a prime. Let $x_1, \ldots, x_c$ be a sequence
in $R$ whose image in $R_{\mathfrak p}$ forms an
$M_{\mathfrak p}$-quasi-regular sequence. Then there exists a
$g \in R$, $g \not \in \mathfrak p$
such that the image of $x_1, \ldots, x_c$ in $R_g$ forms
an $M_g$-quasi-regular sequence.
\end{lemma}

\begin{proof}
Consider the kernel $K$ of the map (\ref{equation-quasi-regular}).
As $M/JM \otimes_{R/J} R/J[X_1, \ldots, X_c]$ is a finite
$R/J[X_1, \ldots, X_c]$-module and as $R/J[X_1, \ldots, X_c]$ is
Noetherian, we see that $K$ is also a finite $R/J[X_1, \ldots, X_c]$-module.
Pick homogeneous generators $k_1, \ldots, k_t \in K$.
By assumption for each $i = 1, \ldots, t$ there exists a $g_i \in R$,
$g_i \not \in \mathfrak p$ such that $g_i k_i = 0$.
Hence $g = g_1 \ldots g_t$ works.
\end{proof}

\begin{lemma}
\label{lemma-truncate-quasi-regular}
Let $R$ be a ring. Let $M$ be an $R$-module.
Let $f_1, \ldots, f_c \in R$ be an $M$-quasi-regular sequence.
For any $i$ the sequence
$\overline{f}_{i + 1}, \ldots, \overline{f}_c$
of $\overline{R} = R/(f_1, \ldots, f_i)$ is an
$\overline{M} = M/(f_1, \ldots, f_i)M$-quasi-regular sequence.
\end{lemma}

\begin{proof}
It suffices to prove this for $i = 1$. Set
$\overline{J} = (\overline{f}_2, \ldots, \overline{f}_c) \subset \overline{R}$.
Then
\begin{align*}
\overline{J}^n\overline{M}/\overline{J}^{n + 1}\overline{M}
& =
(J^nM + f_1M)/(J^{n + 1}M + f_1M) \\
& = J^nM / (J^{n + 1}M + J^nM \cap f_1M).
\end{align*}
Thus, in order to prove the lemma it suffices to show that
$J^{n + 1}M + J^nM \cap f_1M = J^{n + 1}M + f_1J^{n - 1}M$
because that will show that
$\bigoplus_{n \geq 0}
\overline{J}^n\overline{M}/\overline{J}^{n + 1}\overline{M}$
is the quotient of
$\bigoplus_{n \geq 0} J^nM/J^{n + 1}M \cong M/JM[X_1, \ldots, X_c]$
by $X_1$. Actually, we have $J^nM \cap f_1M = f_1J^{n - 1}M$.
Namely, if $m \not \in J^{n - 1}M$, then $f_1m \not \in J^nM$
because $\bigoplus J^nM/J^{n + 1}M$ is the polynomial algebra
$M/J[X_1, \ldots, X_c]$ by assumption.
\end{proof}

\begin{lemma}
\label{lemma-quasi-regular-regular}
Let $(R, \mathfrak m)$ be a local Noetherian ring.
Let $M$ be a nonzero finite $R$-module.
Let $f_1, \ldots, f_c \in \mathfrak m$ be an $M$-quasi-regular sequence.
Then $f_1, \ldots, f_c$ is an $M$-regular sequence.
\end{lemma}

\begin{proof}
Set $J = (f_1, \ldots, f_c)$.
Let us show that $f_1$ is a nonzerodivisor on $M$.
Suppose $x \in M$ is not zero.
By Krull's intersection theorem there exists an integer $r$
such that $x \in J^rM$ but $x \not \in J^{r + 1}M$, see
Lemma \ref{lemma-intersect-powers-ideal-module-zero}.
Then $f_1 x \in J^{r + 1}M$ is an element whose class in
$J^{r + 1}M/J^{r + 2}M$ is nonzero by the assumed structure of
$\bigoplus J^nM/J^{n + 1}M$. Whence $f_1x \not = 0$.

\medskip\noindent
Now we can finish the proof by induction on $c$ using
Lemma \ref{lemma-truncate-quasi-regular}.
\end{proof}

\begin{remark}[Other types of regular sequences]
\label{remark-koszul-regular}
In the paper \cite{Kabele} the author discusses two more
regularity conditions for sequences $x_1, \ldots, x_r$ of elements
of a ring $R$. Namely, we say the sequence is {\it Koszul-regular}
if $H_i(K_{\bullet}(R, x_{\bullet})) = 0$ for $i \geq 1$ where
$K_{\bullet}(R, x_{\bullet})$ is the Koszul complex. The sequence is
called {\it $H_1$-regular} if $H_1(K_{\bullet}(R, x_{\bullet})) = 0$.
One has the implications regular $\Rightarrow$ Koszul-regular $\Rightarrow$
$H_1$-regular $\Rightarrow$ quasi-regular. By examples the author shows that
these implications cannot be reversed in general even if $R$ is a
(non-Noetherian) local ring and the sequence generates the maximal ideal of
$R$. We introduce these notions in more detail in
More on Algebra, Section \ref{more-algebra-section-koszul-regular}.
\end{remark}

\begin{remark}
\label{remark-join-quasi-regular-sequences}
Let $k$ be a field. Consider the ring
$$
A = k[x, y, w, z_0, z_1, z_2, \ldots]/
(y^2z_0 - wx, z_0 - yz_1, z_1 - yz_2, \ldots)
$$
In this ring $x$ is a nonzerodivisor and the image of $y$ in
$A/xA$ gives a quasi-regular sequence. But it is not true that
$x, y$ is a quasi-regular sequence in $A$ because $(x, y)/(x, y)^2$
isn't free of rank two over $A/(x, y)$ due to the fact that
$wx = 0$ in $(x, y)/(x, y)^2$ but $w$ isn't zero in $A/(x, y)$.
Hence the analogue of
Lemma \ref{lemma-join-regular-sequences}
does not hold for quasi-regular sequences.
\end{remark}

\begin{lemma}
\label{lemma-quasi-regular-on-quotient}
Let $R$ be a ring. Let $J = (f_1, \ldots, f_r)$ be an ideal of $R$.
Let $M$ be an $R$-module. Set $\overline{R} = R/\bigcap_{n \geq 0} J^n$,
$\overline{M} = M/\bigcap_{n \geq 0} J^nM$, and denote
$\overline{f}_i$ the image of $f_i$ in $\overline{R}$.
Then $f_1, \ldots, f_r$ is $M$-quasi-regular if and only if
$\overline{f}_1, \ldots, \overline{f}_r$ is $\overline{M}$-quasi-regular.
\end{lemma}

\begin{proof}
This is true because
$J^nM/J^{n + 1}M \cong
\overline{J}^n\overline{M}/\overline{J}^{n + 1}\overline{M}$.
\end{proof}





